\chapter{Conclusioni e sviluppi futuri}
\label{chap:conclusioni}

\section{Sintesi del lavoro}

In questo lavoro ho confrontato teste classiche e ibride
classico-quantistiche per la classificazione morfologica delle galassie.
Il dataset bilanciato di \(5\,628\) immagini ha permesso di definire due task e di
esaminare il ruolo della disponibilità dei dati attraverso subset
stratificati del training. La matrice finale comprende \(640\) run,
organizzate su dieci ripetizioni di holdout stratificato.

La rappresentazione condivisa di ResNet18 ha consentito di
concentrare l'analisi sulle teste. Le baseline B0 e B1 hanno
fornito riferimenti classici; B2 ha controllato la compressione
a quattro o otto componenti; Q e QF hanno distinto l'inserimento
del circuito dall'addestramento dei suoi parametri. Prestazioni,
incertezza, errori per classe e tempi sono stati esaminati
congiuntamente.

Sul Full binario le architetture raggiungono Macro-F1 prossime
a \(0{,}96\), senza evidenza di un vantaggio di Q rispetto a B2.
Sul Full a tre classi Q presenta invece una perdita media di
\(0{,}73\) pp a quattro qubit e \(1{,}67\) pp a otto,
significativa secondo i test e la correzione adottati. I regimi
con pochi dati non mostrano il beneficio quantistico ipotizzato;
l'analisi merger rende visibile una parte importante delle
difficoltà, soprattutto a otto qubit.

Il contributo dell'addestramento circuitale risulta contenuto
sul Full e dipendente dal regime nei subset. L'aumento dei qubit
non determina miglioramenti sistematici. I tempi dei modelli
quantistici simulati sono invece sensibilmente maggiori di quelli
del controllo classico.
Le conclusioni restano circoscritte a dataset, rappresentazione,
circuiti e protocollo considerati.

\section{Contributi e valore del confronto}

Il risultato del lavoro è un confronto riproducibile su dati
astronomici. Immagini, partizioni, configurazioni e predizioni
rimangono collegate negli artefatti, consentendo di risalire
dalle metriche aggregate alle singole esecuzioni.

La combinazione di controlli classici, circuito congelato e
regimi di dati rende il risultato più informativo di una sola
misura di accuratezza. Permette di distinguere ciò che il modello
ibrido riesce a fare da ciò che apporta rispetto alle alternative.
Nelle condizioni esaminate, il circuito non offre un miglioramento
predittivo che compensi il maggiore costo osservato.

Il protocollo costituisce inoltre una base riutilizzabile per
valutare modifiche successive. Nuove rappresentazioni o circuiti
possono essere confrontati mantenendo criteri espliciti di
appaiamento e una distinzione tra analisi confermative ed
esplorative.

\section{Sviluppi futuri}
\label{sec:concl-futuro}

\subsection{Validazione esterna e robustezza del protocollo}

La priorità è verificare la generalizzazione su un campione
indipendente, possibilmente con condizioni osservative o
prevalenze diverse. Partizioni per regione di cielo o sistema
fisico aiuterebbero a controllare correlazioni non eliminate
dalla sola separazione degli identificativi. Analisi di
sensibilità alle soglie di etichettatura potrebbero inoltre
chiarire quanto il confronto dipenda dalla selezione di
esempi più o meno ambigui.

Sul piano statistico, ulteriori valutazioni dovrebbero
considerare esplicitamente la dipendenza tra split e definire
in anticipo famiglie di confronti e differenze di interesse
pratico. Se l'obiettivo fosse dimostrare che due modelli
hanno prestazioni sufficientemente vicine, servirebbe un
protocollo di equivalenza o non inferiorità, non la sola
assenza di significatività.

\subsection{Rappresentazione e ottimizzazione}

Una seconda direzione riguarda l'adeguatezza delle feature
al dominio galattico. Backbone preaddestrati su immagini
astronomiche o un fine-tuning controllato potrebbero cambiare
la qualità della rappresentazione disponibile alle teste.
Ogni modifica dovrebbe essere applicata anche ai controlli
classici per non attribuire al circuito un beneficio dovuto
all'estrattore.

È inoltre utile confrontare strategie di ottimizzazione con
budget di ricerca espliciti e comparabili, distinguendo
numero di epoche, aggiornamenti e tempo disponibile. Questo
permetterebbe di verificare se gli iperparametri condivisi
abbiano favorito alcune teste e se una diversa ottimizzazione
cambi il divario osservato.

\subsection{Ablazioni circuitali e diagnostiche mirate}

Le modifiche al circuito dovrebbero rispondere a ipotesi
specifiche: cambiare profondità, collegamenti o encoding
separatamente è più informativo che aumentare simultaneamente
tutte le componenti. Il data re-uploading, non usato nel
circuito principale, rappresenta una possibile estensione
da confrontare con controlli classici adeguati.

Per comprendere il ruolo dell'addestramento quantistico
occorrerebbe registrare separatamente gradienti e variazioni
dei soli parametri circuitali. Esperimenti con circuito
senza entangling gate e analisi della sensibilità agli
ingressi potrebbero verificare quale parte della trasformazione
sia effettivamente utile. Si tratterebbe di nuovi esperimenti,
non di spiegazioni già dimostrate dai risultati disponibili.

\subsection{Costo e condizioni di esecuzione realistiche}

Un benchmark su hardware costante, con misure end-to-end
e separazione tra estrazione delle feature, training e
inferenza, renderebbe più completa l'analisi dei costi.
L'introduzione di shot finiti e rumore consentirebbe poi
di studiare la sensibilità delle predizioni a condizioni
più vicine all'esecuzione fisica.

Un'eventuale sperimentazione su QPU richiederebbe di misurare
anche preparazione, compilazione, numero di valutazioni e
costi di accesso, evitando di trasferire direttamente i
tempi del simulatore all'hardware quantistico. Il criterio
resterebbe il beneficio ottenuto rispetto a risorse e
prestazioni dei controlli classici.

\section{Considerazione conclusiva}

Gli esperimenti mostrano che piccoli circuiti quantistici possono
essere integrati nella pipeline e raggiungere buone prestazioni,
ma che questo non basta a renderli preferibili al controllo classico.
Con le feature e le condizioni adottate, B2 offre un compromesso
più favorevole tra prestazioni e tempo. Il risultato riguarda le
architetture valutate e non esclude benefici in altri contesti.
Gli sviluppi successivi dovranno verificare quali modifiche alla
rappresentazione, al circuito o all'ottimizzazione cambino questo
confronto, mantenendo controlli e criteri di valutazione espliciti.
